%% Q1 ctx
\textbf{Directions (Q.~Nos.~1 to 5):} Read the para:\par\smallskip
Some matters do not decompose. Some matters take longer time to decompose. All plants decompose. All creepers are plants. Plastics decompose in longer time. Animals decompose in longer time.\par\smallskip
\textbf{Which of the following statement(s) \emph{strictly} do not flow from above given statements?} (Q.~1--5)

%% Q1 q
Which of the following statement(s) strictly do not flow from the above statements?

%% Q1
Nowhere does the paragraph say that plants, animals or ``everything decomposable'' are \emph{matter}. It only says \emph{some} matters decompose. So (A), (B) and (C) all fail to follow strictly: \textbf{all of the above}.
\tip{``Strictly follows'' means it must be deducible from the given statements alone, without outside knowledge.}

%% Q2 q
Which of the following statement(s) strictly do not flow from the above statements?

%% Q2
``All creepers are plants'' and ``All plants decompose'' give ``creepers decompose'', which follows. But nothing says plants (or creepers) are \emph{matter}, so \textbf{``creepers are matters'' does not follow}.

%% Q3 q
Which of the following statement(s) strictly do not flow from the above statements?

%% Q3
``Plastics decompose in longer time'' does not tell us plastics are \emph{matter}, so \textbf{``plastics are matters''} does not strictly follow (answer as per the key). ``Creepers decompose'' does follow.

%% Q4 q
Which of the following statement(s) strictly do not flow from the above statements?

%% Q4
``Some matters take longer time'' does not mean \emph{whatever} takes longer is matter (A); ``matters are creepers'' (B) and ``animals are matters'' (C) are not stated either. \textbf{All of the above} do not follow.

%% Q5 q
Which of the following statement(s) strictly do not flow from the above statements?

%% Q5
``Plastics take time to decompose'' (B) is stated. ``Some matters decompose, others do not'' (C) follows from the first two sentences. \textbf{``Some matters are plants''} (A) does not follow, since plants are never linked to matter.

%% Q6 ctx
\textbf{Directions (Q.~Nos.~6 to 10):} The following is the data (in crores) for educational expenditure (Ed.\ Ex.) for the various plan periods, population (Pop) and Price Index (PI) [Base 1973 $=$ 100]. This is purely hypothetical.
\begin{center}\small
\begin{tabular}{lcccc}\toprule \textbf{Items} & \textbf{VI} & \textbf{VII} & \textbf{VIII} & \textbf{IX}\\\midrule
Ed.\ Ex. & 17.3 & 28.5 & 37.6 & 38.5\\ Pop. & 50 & 55 & 61 & 70\\ PI & 123 & 144 & 189 & 210\\\bottomrule
\end{tabular}
\end{center}
Worked values used below:
\begin{center}\small
\begin{tabular}{lcccc}\toprule & \textbf{VI} & \textbf{VII} & \textbf{VIII} & \textbf{IX}\\\midrule
Increase in Ed.\ Ex.\ over previous plan & -- & 64.7\% & 31.9\% & 2.4\%\\
Per capita Ed.\ Ex.\ (Ed.Ex./Pop) & 0.346 & 0.518 & 0.616 & 0.550\\
Per capita at constant prices ($\times100/$PI) & 0.281 & 0.360 & 0.326 & 0.262\\
Population increase over previous plan & -- & 5 (10\%) & 6 (10.9\%) & 9 (14.8\%)\\\bottomrule
\end{tabular}
\end{center}

%% Q6 q
In which plan period the total increase (\%) in Educational expenditure was maximum?

%% Q6
From VI to VII: $\dfrac{28.5-17.3}{17.3}\times100=64.7\%$, far higher than VIII (31.9\%) or IX (2.4\%). The answer is \textbf{VII}.

%% Q7 note
From the table the maximum per capita expenditure is in Plan VIII (0.616), not IX (0.550). The official key gives (D).

%% Q7
Per capita expenditure $=$ Ed.\ Ex.\ $\div$ Pop: VI 0.346, VII 0.518, \textbf{VIII 0.616}, IX 0.550. By the data, Plan VIII is highest; the key's answer is IX. Mark (D) as per the official key.

%% Q8 note
By the data the smallest increase is in Plan VII (5 crore, 10\%). The official key gives (D).

%% Q8
Population increases: VII $+5$, VIII $+6$, IX $+9$. The minimum is VII; the key gives IX. Mark (D) as per the official key.

%% Q9
Per capita expenditure corrected for price (PI $=100$): $\dfrac{\text{Ed.Ex.}/\text{Pop}}{\text{PI}}\times100$: VI 0.281, \textbf{VII 0.360}, VIII 0.326, IX 0.262. The maximum is \textbf{VII}.

%% Q10 note
The table does not give the total plan allocation, so this cannot be computed from the data. The key gives (C) VIII.

%% Q10
Expenditure as a percentage of total allocation needs the total outlay of each plan, which is not supplied. The official key answer is \textbf{VIII}.

%% Q11
Research produces new knowledge and solutions, so it is an investment, not waste. The appropriate reaction is \textbf{``I disagree with the statement''}.

%% Q12
The most essential attribute is \textbf{adequate interest in and aptitude for research}, which covers both motivation and ability.

%% Q13
A mature researcher keeps an \textbf{open mind about the results}. The hypothesis is tested, not defended, and the findings are accepted whichever way they go.

%% Q14
Per the key, the purpose of research is \textbf{to construct reality}: to build an understanding of reality through systematic inquiry.

%% Q15
\textbf{Observation schedules} record actual behaviour directly, so they give the most valid data. Attitude scales, rating scales and opinionnaires depend on self-report or judgement.

%% Q16
Accepting no absolute truths while appreciating logical alternatives is \textbf{positive relativism}.

%% Q17
``X includes Y'' ($Y\subseteq X$) and ``Y is included in Z'' ($Y\subseteq Z$) tell us nothing about how X and Z relate. Options (A)--(C) could each be true in some case, so \textbf{all are possible}.

%% Q18
The first rational step is to \textbf{check whether the adjacent houses have light}. This tells you whether it is a general power cut or a local fault before you act.

%% Q19
A teacher teaches punctuality best by \textbf{being punctual himself or herself}: example is better than precept.

%% Q20
Per the key, many join teaching because \textbf{teachers serve the entire society}: it is a noble, service-oriented profession.

%% Q21
A good teacher needs subject mastery, teaching competence and clarity of thought and expression: \textbf{all of the above}.

%% Q22
Character is built most effectively by the \textbf{living example of a good teacher} (values are caught, not taught).

%% Q23
Per the key, the teacher should \textbf{tell the student to be in time}: a firm, respectful correction without humiliation or exclusion.

%% Q24
\textbf{Humiliating students} is the most undesirable behaviour. It damages self-esteem and the learning relationship.

%% Q25
Re-directing the question to another student keeps the class engaged and lets peers correct the error without embarrassing the first student.

%% Q26 ctx
\textbf{Direction (Q.~Nos.~26 to 30):} Read the following text and answer questions 26 to 30.\par\smallskip
Television was, towards the end of the war in Vietnam, the main source of information for most Americans, and the images on screen, in black and white, changed war reporting forever. Morley Safer's expose of American soldiers burning a village had a huge impact on viewers at home, and added to the military's distrust of reporters. In Vietnam, correspondents were made honorary majors (World War II reporters had the momentary rank of captain) and could travel anywhere on military transport if space were available. The only condition: that they could not betray troop movements, a request that was respected. The press battles were over misleading briefings which ran contrary to what the reporters could see for themselves in the field.

To this day, some military men glumly insist that television and the press lost the war in Vietnam by demoralizing the home front. The U.S.\ military would never again allow such access to its operations as it did in Vietnam.

The tension between the military and the press will never cease, because both need each other but cannot grant the other what it really wants.

%% Q26
Morley Safer's televised exposé of soldiers burning a village ``added to the military's distrust of reporters'', i.e.\ \textbf{the war images telecast by American television companies}.

%% Q27
Vietnam correspondents were made \textbf{honorary majors} (a higher rank than WWII reporters, who were captains) and could travel on military transport: a higher rank and more facilities.

%% Q28
``The press battles were over misleading briefings which ran contrary to what the reporters could see for themselves'', i.e.\ the \textbf{difference between the military's version and the reporters' version}.

%% Q29
Some military men insist that ``television and the press lost the war in Vietnam by \textbf{demoralizing the home front}'', which the key paraphrases as the mass media turning people against the war.

%% Q30
The passage says both ``need each other but cannot grant the other what it really wants''. The key gives (D): the press pursues truth in all situations, which the military cannot always permit.

%% Q31
Anger, annoyance and irritation are negative emotions of displeasure. \textbf{Smile} is a positive expression, so it does not belong.

%% Q32
Gnashing teeth, knitted brows and creased foreheads express anger or worry. \textbf{Smiling eyes} express happiness.

%% Q33
A bird's body is covered with \textbf{feathers}; a cat's body is covered with \textbf{fur}.

%% Q34
Lead is a typical example of something \emph{heavy}; a \textbf{breeze} is a typical example of something \emph{cool} (per the key).

%% Q35
Air animates (gives life); vitamins give \textbf{health}.

%% Q36
Total age of 32 students $=320$. With the teacher, 33 persons average 11, so the total is $363$. Teacher $=363-320=\mathbf{43}$ years.

%% Q37
Copper fraction $=\dfrac{0.5}{9.5+0.5}=\dfrac{1}{20}$. In an 8 g ring, copper $=\dfrac{8}{20}=\mathbf{0.4\ g}$.

%% Q38
Let $p$ = peacocks and $a$ = animals: $p+a=35$ and $2p+4a=98$, so $2a=98-70=28$, $a=14$ and $p=\mathbf{21}$.

%% Q39
An object looks green because it \textbf{reflects green light} and absorbs the other colours.

%% Q40
FRIEND $\rightarrow$ IOKCQA: the shifts are $+3, -3, +2, -2, +3, -3$.
SPIKER: S$+3=$V, P$-3=$M, I$+2=$K, K$-2=$I, E$+3=$H, R$-3=$O, giving \textbf{VMKIHO}.

%% Q41
``Quality of research depends on the researcher's knowledge of methodology'' is a testable proposition about a relationship, i.e.\ a \textbf{hypothesis}.

%% Q42
``Teaching is an art'' is a value judgement: an \textbf{opinion}.

%% Q43
``Physical growth of a learner is age-related'' is an established \textbf{fact}.

%% Q44
Per the key, ``manipulation of independent variables is required in experimentation'' is treated as an \textbf{assumption}, a premise on which experimental research rests.

%% Q45
``Protein is essential for reconstruction of cells'' is a scientific \textbf{fact}.

%% Q46
The sequence is 1 followed by prime numbers: 1, 3, 5, 7, 11, 13, \textbf{17}.

%% Q47 note
The logical next term is $11\times13=143$, which is not offered. The official key gives (C) 141.

%% Q47
Each term is a product of consecutive odd primes: $1\times3=3$, $3\times5=15$, $5\times7=35$, $7\times11=77$, so next $11\times13=143$. The closest option, and the key's answer, is \textbf{141}.

%% Q48
The differences increase by 2 each time: $+2, +4, +6, +8, +10, +12$, so next $+14$: $45+14=\mathbf{59}$.

%% Q49
Two alternating series: $4, 5, 6$ (increasing) and $2, 1, \mathbf{0}$ (decreasing).

%% Q50 note
$5\times1221=6105$ is not among the options; the official key gives (B) 6006.

%% Q50
The terms are $1221\times1, \times2, \times3, \times4$. The next would be $1221\times5=6105$. The key's intended pattern is palindromic numbers ($a\,2a\,2a\,a$), and it gives \textbf{6006}.
